\documentclass[a4paper,11pt]{article}
\usepackage{exam-style-842}
\examinehead{842 信息化概论}{模拟试卷（二）}
\begin{document}

\examcover{842 信息化概论}{模拟试卷（二）}{信息化概论（依据《信息系统导论》教材）}{150 分}{180 分钟}{本卷卷面为名词解释 30 分、简答题 60 分、论述题 60 分。}

\parttitle{一、名词解释}
\qtypename{名词解释}{6 题，每题 5 分，共 30 分；要求写出概念内涵，必要时举例}
\q 数字化\anslines{2.0cm}
\q 虚拟化技术\anslines{2.0cm}
\q 数据挖掘\anslines{2.0cm}
\q 移动电子商务\anslines{2.0cm}
\q 农业专家系统\anslines{2.0cm}
\q 3S 技术\anslines{2.0cm}
\newpage

\parttitle{二、简答题}
\qtypename{简答题}{6 题，每题 10 分，共 60 分}
\q 简述操作系统的主要功能，并说明为什么操作系统是计算机系统中最基本的系统软件。\anslines{3.2cm}
\q 比较 B/S 架构与 C/S 架构的优缺点，并说明 Web 服务在信息系统集成中所起的作用。\anslines{3.2cm}
\q 简述 MapReduce 的基本原理与执行过程，并说明它在大数据处理中的作用。\anslines{3.2cm}
\newpage
\q 物联网的体系结构通常划分为哪几层？说明各层的功能与代表性技术，并举出一个农业应用实例。\anslines{3.2cm}
\q 试述 3S 技术中遥感、地理信息系统、全球定位系统各自的作用及其在农业中的典型应用。\anslines{3.2cm}
\q 简述移动电子商务的主要模式与特点，并说明电子商务与智能物流之间的关系。\anslines{3.2cm}
\newpage

\parttitle{三、论述题}
\qtypename{论述题}{2 题，每题 30 分，共 60 分；要求观点明确、条理清晰、联系我国农业农村实际}
\q 论述物联网技术体系及其在农业中的应用，并分析农业物联网发展中存在的困难与推进思路。\anslines{11cm}
\newpage
\q 结合数字乡村战略，论述信息技术如何赋能乡村振兴，并分析当前农村农业信息化进程中存在的主要障碍及破解对策。\anslines{13cm}

\end{document}
